% concatenated from neurips_2026.tex
\documentclass{article}

\PassOptionsToPackage{numbers, compress}{natbib}
\usepackage[main, final]{neurips_2026}

\usepackage[utf8]{inputenc}
\usepackage[T1]{fontenc}
\usepackage{hyperref}
\usepackage{url}
\usepackage{booktabs}
\usepackage{amsfonts}
\usepackage{amsmath}
\usepackage{amssymb}
\usepackage{amsthm}
\usepackage{nicefrac}
\usepackage{microtype}
\usepackage{xcolor}
\usepackage{array}

\newtheorem{theorem}{Theorem}
\newtheorem{lemma}[theorem]{Lemma}
\newtheorem{proposition}[theorem]{Proposition}
\newtheorem{corollary}[theorem]{Corollary}
\theoremstyle{definition}
\newtheorem{definition}[theorem]{Definition}
\newtheorem{remark}[theorem]{Remark}

\newcommand{\R}{\mathbb{R}}
\newcommand{\Z}{\mathbb{Z}}
\newcommand{\N}{\mathbb{N}}
\newcommand{\Q}{\mathbb{Q}}
\newcommand{\B}{\{0,1\}}
\newcommand{\norm}[1]{\left\lVert#1\right\rVert}
\newcommand{\eps}{\varepsilon}
\newcommand{\nc}{\mathcal{N}}
\newcommand{\Param}{\Theta}
\newcommand{\halt}{\mathsf{halt}}
\newcommand{\emit}{\mathsf{emit}}
\newcommand{\bit}{\mathsf{bit}}

\title{Neural Weight Norm = Kolmogorov Complexity}

\author{%
  Tiberiu Musat \\
  ETH Zürich \\
  \texttt{tiberiu@musat.ai}
}


\begin{document}

\maketitle

\begin{abstract}
  Why does weight decay work? We prove that, in any fixed-precision regime, the smallest weight norm of a looped neural network outputting a binary string equals the Kolmogorov complexity of that string, up to a logarithmic factor. This implies that weight decay induces a prior matching Solomonoff's universal prior, the optimal prior over computable functions, up to a polynomial factor. The result is norm-agnostic: in fixed precision, every weight norm collapses to the non-zero parameter count up to constants, so the same sandwich bound holds for any norm used as a regulariser. The proof has two short reductions: any program for a universal Turing machine can be encoded into neural weights at unit cost per program bit, and any fixed-precision network can be described by enumerating its non-zero parameters with logarithmic addressing overhead. Both bounds are tight up to constants, with the logarithmic factor realised by permutation encodings: a network whose parameters encode a permutation produces a string whose Kolmogorov complexity is the non-zero parameter count times its logarithm. The fixed-precision assumption is essential: with infinite precision, neural networks can encode non-computable functions and the weight norm loses its relevance.
\end{abstract}

\section{Introduction}\label{sec:intro}

Why does weight decay \citep{krogh1991simple} work? It is the most universally
applied regulariser in modern deep learning. Removing it consistently degrades
held-out performance \citep{andriushchenko2024why}; including it, networks
generalise far beyond what classical capacity-based theory predicts. Yet
capacity-based learning theory, which characterises function classes, cannot
tell a network trained with weight decay apart from one trained without: the
function class is the same, the VC and Rademacher complexities are the same
\citep{zhang2017understanding,nagarajan2019uniform}, and the same architectures
that generalise on real data can be trained to memorise random labels
\citep{zhang2017understanding}. Whatever explains weight decay, it is not about the size of the function class.

A natural alternative is description-length theory: prefer hypotheses with
short descriptions \citep{rissanen1978modeling,grunwald2007mdl}. Its sharpest
form is \emph{Kolmogorov complexity} $K(s)$
\citep{kolmogorov1965three,solomonoff1964formal,chaitin1969length}: the length
of the shortest program $p$ that outputs a string $s$ on a fixed universal
Turing machine $U$. This can be used to construct a \emph{universal prior} $M(s) \propto
  \sum_{p\,:\,U(p)=s} 2^{-|p|}$, dominated in its sum by $2^{-K(s)}$, that
yields Bayesian predictions asymptotically beating every computable predictor
\citep{solomonoff1964formal,hutter2005universal,legg2007universal}. The catch
has always been computability: $M$ is incomputable, so the theory has lived as
an asymptotic ideal rather than as a practical inductive bias one could
actually train under.

This paper closes the gap between the asymptotic ideal and the practical
regulariser. We prove that, in any fixed-precision regime, the smallest weight
norm of a Turing-complete looped neural network outputting a string $s$ equals
the Kolmogorov complexity $K(s)$, up to a logarithmic factor. The same holds
for any $L_p$ norm raised to its corresponding power: in fixed precision, every
weight norm controls the same underlying quantity, the count of non-zero
parameters.

\newpage
\begin{quote}
  \textbf{Main Theorem (informal).} \emph{Let $\nc(s)$ be the minimum number of non-zero parameters in a fixed-precision looped neural network that halts and outputs $s$, and let $K(s)$ be the Kolmogorov complexity of $s$. Then
    \[
      \nc(s) \;\le\; K(s) \;\le\; \nc(s) \log \nc(s),
    \]
    up to constants and with both ends tight. The same sandwich transfers to every
    $L_p$ weight norm, since $\norm{\theta}_p^p = \Theta(\nc(\theta))$ in fixed
    precision.}
\end{quote}

\paragraph{The corollary.} The weight decay penalty term is the negative log of a Gaussian prior. Up to
the logarithmic factor of our bound, the induced prior on the network's output
is Solomonoff's universal prior. The most reliable regulariser in modern deep
learning is, asymptotically, the prior an idealised computationally-unbounded
Bayesian agent would write down as optimal over the space of computable
hypotheses.

\paragraph{Why the logarithmic factor is fundamental.} The log factor is not an artefact of the proof; the network can actually
exploit it. Each non-zero parameter buys $\Theta(\log W)$ bits of address space
(where $W$ is the total non-zero count), and there are families of strings, the
permutations being the canonical example, on which the network cashes out
exactly that budget. For example, it is possible to encode a permutation $\pi :
  [N] \to [N]$ as $N$ ternary edges of a sparse permutation matrix and stream the
matrix row by row. This way, the network uses $\Theta(N)$ ternary parameters
and emits a string of Kolmogorov complexity $\log_2 N! = \Theta(N \log N)$.
Thus, the bounds are two-sided \emph{and} tight.

\paragraph{The role of fixed precision.} Fixed precision is the only neural regime in which the weight-norm-versus-K
question has a non-trivial answer. Real-valued weights are super-Turing
\citep{siegelmann1995computational}: a finite-norm network decides
non-recursively-enumerable languages, so $K(s) = \infty$ while $\norm{\theta}$
is finite. Restricting to rational weights restores Turing-completeness but
does not help: a single rational $p/q$ has bounded magnitude with arbitrarily
many bits of information in its numerator and denominator, so $\norm{\theta}$
still fails to bound $K$ \citep{balcazar1997computational}. Only fixed
precision makes each weight a $O(1)$-bit object, locking magnitude and
description length together so that the question even has content. Modern
hardware uses fp16, bf16, int8, or int4 arithmetic and lives entirely in this
regime. \S\ref{sec:bg:precision} develops the argument.

\paragraph{What is new.} Each ingredient of our argument is well known. Looped transformers are
Turing-complete at constant precision
\citep{li2025constant,li2026efficient,giannou2023looped,perez2021attention,bhattamishra2020computational};
sparse parameter vectors admit short descriptions; the universal prior
dominates all computable predictors. What was missing is a \emph{quantitative
  two-sided} bound between an $L_p$ weight norm of a fixed-precision network and
the Kolmogorov complexity of its emitted string. The closest prior work all
gets near this picture but differs along at least one axis:
\citet{schmidhuber1997discovering} measured complexity by an \emph{external}
program emitting the weight matrix; \citet{jacot2025deep} sandwiched a weighted
L1 ResNet norm against binary circuit gate count for \emph{real-valued}
networks; \citet{shaw2026bridging} bridged transformer MDL to Kolmogorov
complexity through a variational Gaussian-mixture posterior, not via an $L_p$
norm. Section~\ref{sec:related} compares along four axes in a single table.

\paragraph{Outline.} \S\ref{sec:bg} reviews Kolmogorov complexity, Solomonoff induction, and Turing-complete neural networks. \S\ref{sec:setup} fixes the looped network model, the fixed-precision parameter space, and the neural complexity $\nc(s)$. \S\ref{sec:main} states and proves the sandwich bound, including the permutation tightness witness. \S\ref{sec:implications} derives the Solomonoff corollary and an MDL-style generalisation bound. \S\ref{sec:related} contrasts with prior work via a comparison table. \S\ref{sec:discussion} closes with limitations and conjectures. Detailed proofs and extensions are in the appendices.

\section{Background}\label{sec:bg}

\subsection{Kolmogorov complexity and the universal prior}

Fix a universal self-delimiting Turing machine $U$, so that $\{p : U(p)\text{
    halts}\}$ is a prefix-free code over $\B$ \citep{li2008introduction}. The
\emph{(prefix) Kolmogorov complexity} of $s \in \B^*$ is
\[
  K(s) \;=\; \min \bigl\{ \, |p| \,:\, p \in \B^*,\; U(p) = s \,\bigr\}.
\]
The choice of $U$ affects $K$ only by an additive constant; we treat $K$ as
defined up to $O(1)$. The (plain) Kolmogorov complexity $C(s)$ differs from
$K(s)$ by at most an additive logarithmic term, and all our asymptotic
statements hold under either definition.

Solomonoff's \emph{universal prior} is $M(s) \propto \sum_{p \,:\, U(p)=s}
  2^{-|p|}$, dominated in its sum by $2^{-K(s)}$, so $-\log M(s) = K(s) + O(1)$.
Bayesian prediction with $M$ asymptotically dominates every
lower-semicomputable predictor; it is, in a precise sense, the optimal prior an
unbounded reasoner can write down
\citep{solomonoff1964formal,hutter2005universal,legg2007universal}. $K$ is
upper-semicomputable but not computable.

\subsection{Looped neural networks and Turing completeness}

We work with \emph{looped} networks: a single feedforward block $f_\theta$
iterated until a designated halt channel fires. This is the standard formalism
for studying neural networks as computers
\citep{dehghani2018universal,giannou2023looped} and captures the recurrent
depth allocation used by Universal Transformers, deep equilibrium models, and
chain-of-thought reasoning.

A line of work establishes that modern neural architectures are Turing complete
\citep{perez2019turing,perez2021attention,bhattamishra2020computational,Roberts2024decoder,qiu2025prompt}.
Of particular relevance, \citet{li2025constant,li2026efficient} prove Turing
completeness with \emph{constant bit-precision} arithmetic, allocating an
unbounded tape via context length, and \citet{giannou2023looped} give explicit
constant-width constructions that simulate arbitrary register machines. We use
these to instantiate a \emph{universal looped network} $T_U$, of constant size
in fixed precision, that on input $p$ simulates $U(p)$. Concrete instantiations
are in any of the cited works; Appendix~\ref{app:model} fixes one.

\subsection{Fixed precision is the only regime where the question makes sense}\label{sec:bg:precision}

The fixed-precision condition is not a technical convenience. For weight norm
to bound Kolmogorov complexity at all, the $L_p$ norm ball must contain only
finitely many representable weights. Every infinite-precision regime fails this
test, and our bound is correspondingly vacuous in each.

\emph{Real weights are super-Turing.}
\citet{siegelmann1995computational} proved that recurrent networks with
real-valued weights of bounded magnitude decide all languages in exponential
time, including non-recursively-enumerable ones. For such networks
$K(s) = \infty$ on the emitted language while $\norm{\theta}$ is finite, so
$K(s) \le f(\norm{\theta})$ is vacuous for any $f$.

\emph{Rational or computable-real weights still break the bound.} Restricting
to rational weights recovers Turing-completeness without super-Turingness
\citep{siegelmann1995computational}, but the K-complexity bound still fails
for an information-theoretic reason: the ball $\{w \in \Q : |w| \le B\}$ is
countably infinite for any $B$, so $K(w)$ is unbounded over the ball. A single
rational $p/q$ with $|p/q| \le B$ carries arbitrarily many bits of information
via its numerator and denominator; the Siegelmann-Sontag construction
witnesses exactly this, encoding a universal Turing machine into a few
bounded-magnitude rationals. \citet{balcazar1997computational} place rational,
computable-real, and real weight classes in a strict hierarchy indexed by the
Kolmogorov complexity of the weights themselves.

\emph{Fixed precision is the unique resolution.} In any fixed-precision regime,
each non-zero weight is drawn from a finite alphabet of $O(1)$ values; the
$L_p$ ball $\{\theta : \norm{\theta}_p \le B\}$ becomes a finite set whose log
size bounds the Kolmogorov complexity of every weight vector inside it.
Magnitude and description length scale together, so $\norm{\theta}_p$ is the
description length of $\theta$ up to constants. This is precisely what makes
Theorem~\ref{thm:main} possible. Fixed precision is also the regime in which
deep learning actually runs (fp16, bf16, int8, int4), so the bound applies
wherever neural networks are trained or deployed on real hardware. The
fixed-precision assumption is not a simplifying choice but the substantive
content of the result: looped fixed-precision is the only neural model in which
the weight-norm-vs-Kolmogorov question has a non-trivial answer.

\section{The Looped Neural Network Model}\label{sec:setup}

\subsection{Streaming output}\label{sec:setup:model}

\begin{definition}[Looped neural network]\label{def:lnn}
  A \emph{looped neural network} is a pair $(\theta, n)$ where $\theta$ is the parameter vector of an $L$-layer feedforward network $f_\theta : \R^n \to \R^n$ with ReLU activations between layers, and $n \in \N$ is the state dimension. We designate three coordinates of $\R^n$: coordinate $1$ is the \emph{halt} channel $\halt$, coordinate $2$ is the \emph{emit} channel $\emit$, and coordinate $3$ is the \emph{bit} channel $\bit$.

  The network is \emph{run} from the canonical initial state $x_0 = \mathbf{0}$
  by the iteration $x_{t+1} = f_\theta(x_t)$ for $t = 0, 1, 2, \ldots$. At each
  iteration $t \ge 1$, in order:
  \begin{enumerate}
    \item If $(x_t)_{\halt} > 0$, halt.
    \item Otherwise, if $(x_t)_{\emit} > 0$, append the bit $\mathbf{1}[(x_t)_{\bit} >
              0]$ to the output.
  \end{enumerate}
  The network \emph{outputs} $s \in \B^*$ if it halts at some iteration $\tau$ with output bits accumulated so far equal to $s$. If the halt channel never fires, the output is undefined.
\end{definition}

The three-channel convention separates control flow (halt, emit) from data
(bit). It lets the network spend an unbounded number of internal-computation
iterations between emitted bits and so produce arbitrary-length outputs. It
mirrors the way an autoregressive transformer streams tokens. A concrete
realisation as a looped transformer is in Appendix~\ref{app:model}.

\subsection{Fixed precision and the $L_p$ collapse}\label{sec:setup:precision}

\begin{definition}[Fixed-precision parameters]\label{def:precision}
  Fix $\delta > 0$ (resolution) and $M < \infty$ (magnitude bound). The fixed-precision parameter space is
  \[
    \Param_{\delta, M} \;=\; \bigl\{ \theta \in (\delta \cdot \Z \cap [-M, M])^d \,:\, d \in \N \bigr\}.
  \]
\end{definition}

The principal cases of interest fit this template: \emph{ternary} ($\delta = M
= 1$), \emph{$b$-bit signed integer} ($\delta = 1$, $M = 2^{b-1} - 1$), and
\emph{$b$-bit signed dyadic} ($\delta = 2^{-b'}$, $M$ a power of two, closest
to fp16/bf16/int8 quantisation). All are Turing-complete
\citep{li2025constant,giannou2023looped}.

\paragraph{The $L_p$ collapse.} The observation that makes our result norm-agnostic is the following: in any
fixed-precision regime, every $L_p$ norm to the $p$th power equals the non-zero
parameter count up to constants.
\begin{equation}\label{eq:lp-collapse}
  \delta^p \cdot \norm{\theta}_0 \;\le\; \norm{\theta}_p^p \;=\; \sum_{i=1}^d |\theta_i|^p \;\le\; M^p \cdot \norm{\theta}_0
  \qquad \text{for all } p \in [1, \infty).
\end{equation}
The lower bound holds because every non-zero $|\theta_i| \ge \delta$; the upper bound because every $|\theta_i| \le M$. For ternary parameters the bounds coincide ($\norm{\theta}_p^p = \norm{\theta}_0$ exactly); for any other fixed precision they coincide up to a multiplicative constant. Regularising with L1, the squared L2, or any $L_p$ norm controls the same underlying quantity: the count of non-zero parameters.

\subsection{Neural complexity}

\begin{definition}[Neural complexity]\label{def:nc}
  For $s \in \B^*$,
  \[
    \nc(s) \;=\; \inf_\theta \bigl\{ \norm{\theta}_0 \,:\, \theta \in \Param_{\delta, M},\; \theta \text{ outputs } s \bigr\}.
  \]
  We treat $(\delta, M)$ as fixed constants of the model and suppress them. Set
  $\nc(s) = \infty$ if no fixed-precision looped network outputs $s$.
\end{definition}

By Turing completeness in fixed precision, $\nc(s) < \infty$ iff $s$ is
computable. By (\ref{eq:lp-collapse}), the $L_p$-norm minimum $\nc_p(s) :=
  \inf_\theta \{ \norm{\theta}_p : \theta \text{ outputs } s \}$ satisfies
$\nc_p(s)^p = \Theta(\nc(s))$, so $\nc$ is the canonical quantity and bounds on
any $L_p$ norm follow as immediate corollaries.

\section{The Sandwich Bound}\label{sec:main}

\begin{theorem}[Main]\label{thm:main}
  Let $s \in \B^*$ be computable. There exist constants $c_U, c_d \ge 0$,
  depending on the universal Turing machine $U$, the architecture, and the
  precision $(\delta, M)$ but not on $s$, such that
  \begin{align}
    \nc(s) & \;\le\; K(s) + c_U, \label{eq:upper}                           \\
    K(s)   & \;\le\; c_d \cdot \nc(s) \log_2 \nc(s) + c_d. \label{eq:lower}
  \end{align}
  The same sandwich transfers to every $L_p$ weight norm via the collapse
  $\norm{\theta}_p^p = \Theta(\nc(\theta))$ (Proposition~\ref{prop:lp}).
\end{theorem}

The proof is two short reductions: one direction encodes any $U$-program into a
network at unit cost per program bit; the other recovers a program from any
fixed-precision network at $O(\log W)$ bits per non-zero parameter. We give the
body proofs in \S\ref{sec:main:reductions}; tightness, including the
permutation witness, in \S\ref{sec:main:tightness}; full proofs in
Appendices~\ref{app:upper} and~\ref{app:lower}.

\subsection{The two reductions}\label{sec:main:reductions}

\begin{lemma}[Programs $\to$ networks]\label{lem:upper}
  For every program $p \in \B^*$ such that $U(p)$ halts, there is a fixed-precision
  looped network whose output equals $U(p)$ and whose non-zero parameter count
  is at most $|p| + c_U$, where $c_U \ge 0$ is a constant independent of $p$.
\end{lemma}

\begin{proof}[Proof]
  The idea: take a universal looped network and pre-load $p$ into its input
  region via $|p|$ scalar weights, one per program bit.

  Let $T_U$ be the looped network of \citet{li2025constant,giannou2023looped}
  that simulates $U$: when a program is placed in a designated input region of
  $T_U$'s state at iteration $1$, $T_U$ halts and emits its output on the
  emit/bit channels. Write $c_U'$ for the non-zero count of $T_U$, a constant
  independent of $p$.

  Since our model fixes $x_0 = \mathbf{0}$, we cannot place $p$ in the initial
  state directly; we instead inject $p$ into the residual stream once, at the
  first iteration, via a gate that fires only at $t = 1$. Augment $T_U$ with two
  pieces:
  \begin{itemize}
    \item An \emph{iteration-1 gate}: a single state coordinate $c$ with $c^{(1)} = 1$
          and $c^{(t)} = 0$ for $t \ge 2$, driven by a fixed two-line recurrence
          (Appendix~\ref{app:upper}). $O(1)$ non-zero parameters.
    \item A \emph{routing layer}: for each $i \in [|p|]$, one ternary weight $\sigma_i
            \in \{-\delta, +\delta\}$ from $c$ to coordinate $3+i$ of the state, with
          $\sigma_i$ positive if $p_i = 1$ and negative otherwise. Exactly $|p|$ non-zero
          parameters.
  \end{itemize}
  At $t = 1$, the routing layer adds $(\sigma_1, \ldots, \sigma_{|p|})$ to
  coordinates $4, \ldots, 3 + |p|$, depositing $p$ into $T_U$'s input region;
  the residual stream then carries $p$ forward unchanged. At $t \ge 2$,
  $c = 0$ and the routing layer contributes nothing, so $T_U$ runs on the
  loaded program. The constant-precision construction of
  \citet{li2025constant} requires the program to be spread across positions
  rather than packed into a single position (each position holds only $O(1)$
  bits at constant width), so each routing weight is best read as a
  position-specific bias on a distinct sequence position; their tape encoding
  needs $\Theta(|p|)$ positions, matching the $\Theta(|p|)$ routing weights
  here.

  \[
    \norm{\theta_p}_0 \;=\; \underbrace{c_U'}_{T_U} + \underbrace{|p|}_{\text{routing}} + \underbrace{O(1)}_{\text{gate}} \;=\; |p| + c_U, \qquad c_U := c_U' + O(1).
  \]
  Choosing $p$ as a shortest $U$-program for $s$ gives $\nc(s) \le K(s) + c_U$.
\end{proof}

\begin{lemma}[Networks $\to$ programs]\label{lem:lower}
  There is a constant $c_d > 0$ such that every fixed-precision looped network
  $\theta \in \Param_{\delta, M}$ with $\norm{\theta}_0 \ge 1$ that outputs
  $s$ satisfies
  \[
    K(s) \;\le\; c_d \cdot \norm{\theta}_0 \log_2 \norm{\theta}_0 + c_d.
  \]
\end{lemma}

\begin{proof}[Proof]
  A fixed-precision network is a finite, discrete object: $W$ non-zero parameters, each a tuple (location, value). We exhibit a self-delimiting encoding of $\theta$ of length $O(W \log W)$ bits, then prepend a constant-size simulator.

  \emph{Address space.} Every non-zero parameter touches at most two neurons, so the $W$ non-zero parameters together touch at most $2W$ distinct neurons. Prune unused neurons and renumber. After pruning, each parameter is specified by a tuple
  \[
    \bigl(\,\text{layer index},\; \text{source neuron},\; \text{target neuron},\; \text{value}\,\bigr) \in [W] \times [2W] \times [2W] \times \mathcal{V},
  \]
  where $\mathcal{V} := (\delta\Z \cap [-M, M]) \setminus \{0\}$ has size $O(1)$.
  The number of layers needed is at most $W$ (collapse empty layers). Each tuple
  takes $\le 3 \log_2 W + O(1)$ bits in fixed-width binary. For $W$ tuples plus
  $O(\log W)$ self-delimited metadata, the total encoding length is $3 W \log_2 W
    + O(W)$ bits.

  \emph{Simulator.} A constant-size program $\Pi$ parses the encoding, reconstructs $\theta$, and simulates the network forward at fixed precision until the halt channel fires, emitting the bits accumulated on the emit/bit channels. Prepending $\Pi$ gives a self-delimited program $\hat p(\theta)$ for $U$ with $U(\hat p(\theta)) = s$ and $|\hat p(\theta)| \le c_d W \log_2 W + c_d$ for $c_d = \max(3 + \eps, |\Pi|)$.

  Hence $K(s) \le |\hat p(\theta)| \le c_d W \log W + c_d$.
  Appendix~\ref{app:lower} gives the explicit encoding.
\end{proof}

\subsection{Both bounds are tight}\label{sec:main:tightness}

The upper bound is met by an explicit construction: Lemma~\ref{lem:upper}
produces, for every computable $s$, a fixed-precision looped network outputting
$s$ with at most $K(s) + c_U$ non-zero parameters. We now show the lower
bound's logarithmic factor is met as well, by exhibiting a family of strings on
which $K(s) = \Theta(\nc(s) \log \nc(s))$.

\paragraph{The permutation example.} Let $\pi : [N] \to [N]$ be a permutation and $P_\pi \in \{0,1\}^{N \times N}$
its permutation matrix. Define $s_\pi \in \B^{N^2}$ to be the row-major
serialisation of $P_\pi$. Two facts: (i) $K(s_\pi) = \log_2 N! + O(\log N) =
  \Theta(N \log N)$ for typical $\pi$, since $\pi \mapsto s_\pi$ is a bijection
and there are $N!$ permutations. (ii) A looped network with $\Theta(N)$ ternary
parameters can output $s_\pi$ by scanning $(i, j) \in [N]^2$ and emitting
$(P_\pi)_{i,j}$ at each step. Of these $\Theta(N)$ parameters, exactly $N$
encode the choice of $\pi$ through their positions (each a $+1$ entry of
$P_\pi$, contributing $\log_2 N$ bits of address-space information); the
remaining $\Theta(N)$ are a $\pi$-independent control circuit shared across all
permutations. So $\nc(s_\pi) = \Theta(N)$, giving
\[
  K(s_\pi) \;=\; \Theta\!\bigl(\nc(s_\pi)\,\log \nc(s_\pi)\bigr),
\]
which saturates (\ref{eq:lower}) up to constants.
Appendix~\ref{app:permutation} gives the explicit construction and shows why a
binary-output variant of the same example would, by contrast, fail to saturate.

\paragraph{Why the log factor is fundamental.} Each non-zero parameter pays $\log W$ bits to specify its location among the
$\Theta(W^2)$ possible (source, target) pairs in a sparse layout, and the
permutation family realises this entire budget. The bound is two-sided
\emph{and} tight for unstructured architectures.

\paragraph{Could architectural structure improve it?} Restricting the per-parameter address space to some $w_{\mathrm{eff}}(W) \ll
  W^2$ would in principle tighten the log factor to $\log w_{\mathrm{eff}}$.
Encoding obstructs this: with $x_0 = \mathbf{0}$ the program must live in the
weights, and natural restrictions sacrifice that capacity.
Translation-invariant convolutions tie all kernel weights at a given offset
across positions, fitting only $O(1)$ bits of program in the kernel;
fixed-width fixed-precision MLPs are finite-state; looped transformers can
encode arbitrary programs via the routing trick of Lemma~\ref{lem:upper}, but
those routing weights themselves use $w_{\mathrm{eff}} = \Theta(W)$.
Identifying a natural Turing-complete architecture with $w_{\mathrm{eff}} \ll
  W$ is an open question.

\section{Implications}\label{sec:implications}

\subsection{The induced output prior matches Solomonoff's}\label{sec:implications:prior}

The L2 weight-decay penalty $\tfrac{\lambda}{2}\norm{\theta}_2^2$ is the
negative log of a Gaussian prior $\pi(\theta) \propto
  \exp(-\tfrac{\lambda}{2}\norm{\theta}_2^2)$ on weights. By
(\ref{eq:lp-collapse}), the Gaussian and the sparsity prior $\pi_0(\theta)
  \propto 2^{-\norm{\theta}_0}$ agree up to a constant factor in the exponent:
$\delta^2 \norm{\theta}_0 \le -\log_2 \pi(\theta) \le M^2 \norm{\theta}_0$.
Equivalently, $\pi(\theta) = \pi_0(\theta)^{\beta(\theta)}$ for some
$\beta(\theta) \in [\delta^2, M^2]$.

\paragraph{What is compared.} We compare the \emph{induced output prior}
\[
  Q(s) \;:=\; \sum_{\theta \,:\, \theta \text{ outputs } s} \pi(\theta)
\]
against Solomonoff's universal prior $M(s)$. $Q(s)$ is the marginal probability
of emitting $s$ under the joint distribution given by $\pi$ on $\theta$
together with the network's deterministic output map. This is a statement about
the prior at the level of outputs, \emph{not} about training: standard SGD with
L2 weight decay performs MAP estimation of $\theta$, not Bayesian
marginalisation; the claim below is that the prior \emph{itself}, viewed at the
level of the network's output, coincides with Solomonoff's up to a logarithmic
factor.

\begin{corollary}[Induced prior matches Solomonoff]\label{cor:prior}
  Let $\pi(\theta) \propto 2^{-\norm{\theta}_2^2}$ on $\Param_{\delta, M}$ (or,
  up to multiplicative constants in $(\delta, M)$, any $L_p$ analogue), and let
  $Q(s) = \sum_{\theta\,:\,\theta\text{ outputs } s} \pi(\theta)$. There exist
  constants $\alpha, \beta > 0$ (depending on $(\delta, M)$ and the architecture
  but \emph{not} on $s$) such that, for every computable $s$,
  \[
    2^{-K(s) - \alpha} \;\le\; Q(s) \;\le\; 2^{-K(s) / (\beta \log K(s))}.
  \]
\end{corollary}

Equivalently, $-\log Q(s) \in [K(s)/(\beta \log K(s)),\; K(s) + \alpha]$ for
every computable $s$, where Solomonoff's prior satisfies $-\log M(s) = K(s) +
  O(1)$ \citep{li2008introduction}. The induced output prior under L2 weight
decay matches Solomonoff's in the exponent up to a logarithmic factor, with
constants uniform in $s$. To our knowledge, this is the first quantitative
two-sided link between a regulariser used in practice and the universal prior.
Proof in Appendix~\ref{app:prior}.

\subsection{An MDL-style generalisation bound}

The encoding $\hat p(\theta)$ of Lemma~\ref{lem:lower} is a prefix-free
injection into $\B^*$ with $|\hat p(\theta)| \le c_d \norm{\theta}_0 \log
  \norm{\theta}_0 + c_d$, so by the Kraft inequality the assignment $\pi(\theta)
  := 2^{-|\hat p(\theta)|}$ is sub-distributional. Plugging $\pi$ into the
standard Occam-bound / PAC-Bayes argument
\citep{arora2018stronger,lotter2022pacbayes} gives the following.

\begin{proposition}[MDL generalisation bound]\label{prop:gen}
  Fix a confidence level $\eta \in (0, 1)$. With probability at least $1 - \eta$
  over a sample $S$ of size $m$, every fixed-precision network $\theta \in
    \Param_{\delta, M}$ satisfies
  \[
    L(\theta) \;\le\; \hat L(\theta) \;+\; \widetilde O\!\left(\sqrt{\frac{c_d \norm{\theta}_2^2 \log \norm{\theta}_2^2 + \log(1/\eta)}{m}}\right),
  \]
  where $L, \hat L$ are population and empirical risk and $\widetilde O$ hides
  factors logarithmic in $m, 1/\eta$, and the precision $(\delta, M)$.
\end{proposition}

The novelty is conceptual rather than quantitative: existing PAC-Bayes
compression bounds \citep{arora2018stronger,lotter2022pacbayes} achieve the
same shape via more elaborate compression schemes; the bound here uses no
compression beyond the trivial one given by the weights themselves, and the
complexity penalty $\norm{\theta}_2^2 \log \norm{\theta}_2^2$ is expressed in
the same units as the universal prior of Corollary~\ref{cor:prior}.

\subsection{Empirical predictions}\label{sec:implications:predictions}

The result suggests testable conjectures about deep learning practice. None of
these is established here; they are predictions the theorem makes that prior
quantisation-aware training and sparsity literature
\citep{hubara2017quantized,han2016deep,hoefler2021sparsity,molchanov2017variational}
positions us to test directly.

\begin{itemize}
  \item \textbf{Simple data benefits more from weight decay.} Tasks whose targets have low Kolmogorov complexity (algorithmic reasoning, regular languages, structured prediction) should gain more from L2 weight decay than tasks whose targets are essentially random.
  \item \textbf{Quantisation strengthens the bias.} For deeply quantised networks (int4, int8), $\norm{\theta}_2^2$ tracks non-zero parameter count exactly \citep{hubara2017quantized,han2016deep}, so quantisation-aware sparse training is a more direct implementation of the induced output prior than full-precision weight decay.
  \item \textbf{Effective complexity scales as $\norm{\theta}_2^2 \log \norm{\theta}_2^2$.} The right effective complexity for predicting generalisation, in the spirit of \citet{arora2018stronger}, should be the log-augmented squared L2 norm, not raw parameter count or unaugmented norm.
  \item \textbf{Looped depth helps low-$K$ targets.} Networks whose computation budget grows with input difficulty (looped transformers, deep equilibrium models, chain-of-thought) should outperform fixed-depth feedforward networks on targets of bounded but variable Kolmogorov complexity, since the looped model spends description length on state rather than on depth.
  \item \textbf{Sparsity-inducing priors converge to the same prior in fixed precision.} L1 weight decay, magnitude pruning \citep{han2016deep}, and variational dropout \citep{molchanov2017variational} all promote sparse weights; by (\ref{eq:lp-collapse}), in fixed precision they target the same underlying $\norm{\theta}_0$ and therefore induce equivalent output priors up to constants.
\end{itemize}

\section{Related Work}\label{sec:related}

We summarise the closest work in Table~\ref{tab:related} and discuss in turn.

\begin{table}[h]
  \caption{Comparison of weight-norm-vs-complexity bounds. ``Two-sided'' indicates whether both directions of a sandwich are proved.}
  \label{tab:related}
  \centering
  \small
  \begin{tabular}{p{0.22\linewidth}p{0.22\linewidth}p{0.18\linewidth}p{0.13\linewidth}p{0.10\linewidth}}
    \toprule
    \textbf{Reference}                                    & \textbf{Quantity bounded}                          & \textbf{Regulariser}               & \textbf{Precision}           & \textbf{Two-sided?}         \\
    \midrule
    \citet{schmidhuber1997discovering}                    & K-complexity of weight matrix                      & Levin search over external program & Any                          & No                          \\
    \citet{hinton1993keeping}                             & Description length of noisy weights                & Variational MDL                    & Noisy                        & No                          \\
    \citet{hochreiter1997flat}                            & Description length                                 & Flat-minima width                  & Limited                      & No                          \\
    \citet{arora2018stronger}, \citet{lotter2022pacbayes} & Generalisation gap                                 & Compression rate                   & Quantised                    & One-sided                   \\
    \citet{jacot2025deep}                                 & Binary circuit gate count of approximated function & Weighted L1 ResNet norm            & Real-valued                  & \textbf{Yes}                \\
    \citet{shaw2026bridging}                              & K-complexity of distribution                       & Variational MDL (Gaussian mixture) & Variational                  & One-sided (asymp.\ optimal) \\
    \midrule
    \textbf{This paper}                                   & \textbf{K-complexity of emitted string}            & \textbf{Any $L_p$ norm}            & \textbf{Any fixed precision} & \textbf{Yes}                \\
    \bottomrule
  \end{tabular}
\end{table}

\paragraph{Schmidhuber 1997: K-complexity of the weights themselves.}
\citet{schmidhuber1997discovering} measured a network's complexity as the length of an \emph{external} program (in a fixed universal language) that emits the weight matrix, then performed Levin search over those programs. The complexity quantity there is the description length of the external program, not a norm of the weights. Our result is in a sense the converse: under fixed precision, the network's own $L_p$ weight norm is itself such a description length, up to a logarithmic factor, with no external program necessary.

\paragraph{Jacot 2025: closest analog, different complexity.}
\citet{jacot2025deep} proved a sandwich between a weighted L1 norm of real-valued ResNets and the binary circuit gate count of the function being $\eps$-approximated, with equivalence ``within a power of 2'' of the optimal circuit. Read in parallel: Jacot bounds gate count (combinational complexity) of a real-valued ResNet $\eps$-approximating a real function; we bound Kolmogorov complexity (Turing complexity) of a fixed-precision looped network exactly emitting a string. The architectures differ (feedforward ResNet vs.\ looped); the precision differs (real-valued vs.\ fixed); the complexity quantity differs (gates vs.\ programs); the regulariser differs (weighted L1 vs.\ any $L_p$). The two results together suggest that a weight-norm-flavoured quantity is the right complexity measure across a range of architectures and idealisations, with the specific complexity it equals determined by the precision and architecture.

\paragraph{Shaw et al.\ 2026: closest in motivation, different in mechanism.}
\citet{shaw2026bridging} construct an asymptotically optimal MDL objective for transformers via a bridge to prefix Turing machines, achieving description-length optimality up to an additive constant. Their description length is the negative log marginal likelihood under a Gaussian-mixture variational posterior over weights. They optimise a learned variational MDL objective and prove asymptotic optimality of that objective; we show that the squared L2 weight norm \emph{itself}, used directly as a regulariser, is already within a logarithmic factor of optimal. The mechanisms are different (variational MDL vs.\ direct weight norm), but the two papers converge on the picture that standard transformer training is closer to optimal MDL than the absence of prior theory would suggest.

\paragraph{Universal induction and meta-learning.}
\citet{solomonoff1964formal,hutter2005universal,legg2007universal} establish the optimality of the universal prior. \citet{grau2024learning} show that meta-trained transformers can amortise Solomonoff induction empirically. None of these gives a quantitative bound between a weight-space regulariser and the universal prior, which is what Corollary~\ref{cor:prior} provides.

\paragraph{MDL, flat minima, compression bounds.}
\citet{hinton1993keeping} and \citet{hochreiter1997flat} relate generalisation to the precision required to express the weights: their description length is bits-to-represent-noisy-weights, not a norm. \citet{arora2018stronger} and \citet{lotter2022pacbayes} compress trained networks to obtain non-vacuous generalisation bounds, bounding the generalisation gap rather than the Kolmogorov complexity of the emitted string. Theorem~\ref{thm:main} unifies these: in fixed precision, the L2 norm directly counts the structural information in the weights, and a generalisation bound is then a one-step consequence (\S\ref{sec:implications}).

\paragraph{Turing completeness at fixed precision.}
The simulation results we depend on are
\citet{perez2019turing,perez2021attention,bhattamishra2020computational,Roberts2024decoder,qiu2025prompt,li2025constant,li2026efficient,giannou2023looped,dehghani2018universal}.
Of these, \citet{li2025constant,li2026efficient} are most relevant, since they
explicitly handle the constant-bit-precision regime our model captures.

\paragraph{Classical complexity-Kolmogorov links.}
The address-space argument underlying our log factor (each non-zero parameter
buys $\log W$ bits of structural information through its position) has a long
prehistory in complexity theory. Boolean circuit Kolmogorov complexity
\citep{allender2001nlnp} studies the analogous question for circuits emitting
strings, with a similar log factor for wire-indexing. Our contribution is to
specialise this principle to the fixed-precision neural setting and to recover
the matching upper bound via a Turing-complete construction.

\paragraph{Quantisation and sparsity.}
A practical literature on quantisation-aware training
\citep{hubara2017quantized,han2016deep} and sparsity-inducing priors
\citep{molchanov2017variational,frankle2019lottery,hoefler2021sparsity}
operationalises the regime where (\ref{eq:lp-collapse}) is tightest. These
methods target $\norm{\theta}_0$ either directly (pruning) or indirectly
(through quantised L2 or Bayesian-sparsity priors); our result identifies the
underlying complexity quantity they implicitly control.

\section{Discussion}\label{sec:discussion}

\paragraph{Limitations.}
The bound is asymptotic and we do not attempt to make the constants $c_U, c_d$
small; they will be large for realistic universal Turing machines, so the
result is conceptual rather than predictive at small scales. Our model is
restricted to looped networks emitting a single string from $x_0 = \mathbf{0}$;
standard supervised-learning settings (feedforward networks, batched inputs,
training losses) require adaptations we do not pursue here, such as loop
unrolling for feedforward and designated state coordinates for inputs and
labels. Corollary~\ref{cor:prior} concerns the induced output prior $Q(s)$, not
what gradient descent finds; closing the gap between minimum-norm networks and
trained networks is the natural next step. The theorem and its predictions are
not validated experimentally. Spectral and operator norms control Lipschitz
constants rather than information content, and their relation to $K(s)$ remains
open.

\paragraph{Two takeaways.} \emph{(1) The result is genuinely about the fixed-precision regime.} It is the discreteness of $\Param_{\delta, M}$ that gives a finite description length, and the analog limit \citep{siegelmann1995computational} shows no such bound can hold without it. \emph{(2) The result is genuinely about any $L_p$ norm.} The proof concerns the count of non-zero parameters, and the L1-vs-L2 distinction collapses in fixed precision (\ref{eq:lp-collapse}). The penalty's behaviour during training depends on $p$; the resulting Kolmogorov-complexity profile does not.

\paragraph{Conjectures.}
\begin{enumerate}
  \item \emph{Weight norm tracks data complexity.} For networks trained to fit a
        dataset $S$ with L2 weight decay, $\norm{\theta}_2^2$ at convergence tracks
        $K(S)$ in expectation, with the ratio bounded above by a constant as model
        width and depth grow.
  \item \emph{Flat minima are downstream.} The flat-minima phenomenon
        \citep{hochreiter1997flat} is a consequence of the weight-norm-as-description-length
        correspondence: minima of low $L_p$ norm are minima of low description length, and
        small description length is precisely the condition under which the loss landscape
        is locally flat in MDL coordinates.
  \item \emph{Effective complexity is log-augmented.} The right effective complexity for
        predicting generalisation scales as $\norm{\theta}_2^2 \log \norm{\theta}_2^2$,
        not as raw parameter count or as unaugmented norm.
\end{enumerate}

\paragraph{Conclusion.} In any fixed-precision regime, the smallest $L_p$ weight norm of a looped
network outputting $s$ equals $K(s)$ up to a logarithmic factor, realised by
permutation encodings. L2 weight decay therefore matches Solomonoff's universal
prior on outputs up to a polynomial factor in $2^{K(s)}$. Modern deep
learning's most reliable regulariser is a tractable proxy for the most powerful
idealised inductive bias known.

\newpage
\bibliographystyle{plainnat}
\bibliography{refs}

%%%%%%%%%%%%%%%%%%%%%%%%%%%%%%%%%%%%%%%%%%%%%%%%%%%%%%%%%%%%

\appendix

\section{A Concrete Looped Neural Network Model}\label{app:model}

We instantiate Definition~\ref{def:lnn} as a ternary-weight looped transformer,
mirroring \citet{giannou2023looped} with the constant-bit-precision
quantisation of \citet{li2025constant}. The model is the same as in those works
modulo the explicit identification of the halt, emit, and bit channels.

\paragraph{Architecture.} A single transformer block applied iteratively, with: a multi-head
self-attention layer with $h$ heads of dimension $d_h$ and ternary
key/query/value projection matrices; a two-layer feedforward network with
hidden width $d_f$, ReLU activation, and ternary weights and biases; and
residual connections around both sub-layers. The state at iteration $t$ is
$X^{(t)} \in \R^{n \times d}$. The first row is the \emph{control row}; its
coordinate 1 is the halt channel, coordinate 2 is the emit channel, coordinate
3 is the bit channel. Initial state $X^{(0)} = \mathbf{0}$.

\paragraph{Ternary precision.} All parameters lie in $\{-1, 0, +1\}$. Activations are clipped to $\{-A,
  \ldots, +A\}$ for a constant $A$ and rounded after each non-linearity, so the
entire computation lives in a finite alphabet. \citet{li2025constant} prove
that this regime supports Turing-complete simulation; we use their construction
as our universal looped network $T_U$.

\paragraph{Adaptation to other precisions.} The construction is invariant under uniform rescaling of weights and clipping
bounds: replacing $(\theta, A)$ with $(\delta \theta, \delta A)$ preserves all
transitions, so the same architecture instantiates
Definition~\ref{def:precision} in any $\Param_{\delta, M}$ with $\delta \le 1
  \le M$.

\section{Detailed Proof of the Upper Bound}\label{app:upper}

We expand the construction of Lemma~\ref{lem:upper}.

\paragraph{The universal looped network $T_U$.} Let $T_U$ be the ternary looped transformer of Appendix~\ref{app:model}, of
total parameter count $d_U$ and non-zero count $c_U' = \norm{T_U}_0$. By
construction \citep{li2025constant}, $T_U$ has a designated \emph{program
  region} consisting of coordinates $\{4, 5, \ldots, 3 + \ell_{\max}\}$ of the
control row for some maximum program length $\ell_{\max}$ that grows with the
simulation. (In the constant-bit-precision construction of
\citet{li2025constant}, $\ell_{\max}$ scales with the input sequence length;
for our purposes we treat $\ell_{\max}$ as effectively unbounded by feeding a
long enough scratch sequence.)

When $T_U$ is run with the bits of a program $p \in \B^m$ initialised at
positions $4, \ldots, 3 + m$ of the control row, $T_U$ halts (control-row
coordinate 1 becomes positive) at some finite iteration $\tau$ and emits the
bits of $U(p)$ on the emit/bit channels of the control row, provided $U(p)$ is
defined.

\paragraph{Program-loading module.} We construct $\theta_p$ by augmenting $T_U$ with two pieces: a constant-size
\emph{iteration-1 indicator} that fires only on the first iteration, and a
\emph{routing layer} containing exactly one ternary parameter per program bit.

\emph{(i) Iteration-1 indicator.} Add two scratch coordinates $c_1, c_2$ to the state, both
initialised to $0$. Their FFN updates implement the recurrences
\[
  c_2^{(t+1)} = \mathrm{clip}_{[0,1]}\!\bigl(c_2^{(t)} + 1\bigr), \qquad
  c_1^{(t+1)} = \mathrm{ReLU}\!\bigl(1 - c_2^{(t)}\bigr).
\]
Tracing from $c_1^{(0)} = c_2^{(0)} = 0$:
\begin{itemize}
  \item $t = 1$: $c_2^{(1)} = 1$, $c_1^{(1)} = \mathrm{ReLU}(1 - 0) = 1$.
  \item $t \ge 2$: $c_2$ stays at $1$, so $c_1^{(t)} = \mathrm{ReLU}(1 - 1) = 0$.
\end{itemize}
So $c_1^{(t)} = \mathbf{1}[t = 1]$. The recurrences use four ternary parameters total
(one self-connection on $c_2$, one bias on $c_2$, one cross-weight $-1$ from
$c_2$ to $c_1$, one bias $+1$ on $c_1$) plus reuse of the standard
activation clipping primitive of \citet{li2025constant}. Cost: $O(1)$
$p$-independent non-zeros.

\emph{(ii) Program-routing layer.} For each $i \in \{1, \ldots, |p|\}$, place a single ternary parameter $\sigma_i \in
  \{-\delta, +\delta\}$ on the weight from $c_1$ to coordinate $3 + i$ of the
control row, with $\sigma_i = +\delta$ if $p_i = 1$ and $\sigma_i = -\delta$ if
$p_i = 0$. The contribution to coordinate $3 + i$ of the residual stream at
iteration $t$ is therefore $\sigma_i \cdot c_1^{(t)} = \sigma_i \cdot \mathbf{1}[t=1]$.

Trace the residual stream at the program region (coordinates $4, \ldots, 3 +
  |p|$):
\[
  X^{(0)}_{3+i} = 0, \qquad X^{(1)}_{3+i} = X^{(0)}_{3+i} + \sigma_i = \sigma_i, \qquad X^{(t)}_{3+i} = X^{(t-1)}_{3+i} + 0 = \sigma_i \quad \text{for } t \ge 2.
\]
The program bits land in the program region at iteration 1 and persist for
$T_U$'s subsequent simulation. Activation thresholding maps $\pm \delta \to \pm
  1$, the ternary encoding of program bits $T_U$ expects.

\emph{Position-spread encoding.} The construction writes $|p|$ ternary
parameters to $|p|$ distinct coordinates of $T_U$'s state. In the
constant-precision transformer of \citet{li2025constant}, these $|p|$
coordinates correspond to $|p|$ distinct sequence positions, one program bit
per position: at constant model width $d$, a single position holds only $O(d)
  = O(1)$ bits, so an arbitrarily long program cannot fit in a single position
and must be spread across $\Theta(|p|)$ positions of the input sequence. Each
$\sigma_i$ thus acts as a position-specific bias, breaking translation
invariance via a parameter dedicated to position $i$.

\emph{Cost.} The routing layer uses exactly $|p|$ ternary parameters, one per
program bit. No other per-bit overhead.

\paragraph{Counting non-zero parameters.}
\[
  \norm{\theta_p}_0 \;=\; \underbrace{c_U'}_{T_U} \;+\; \underbrace{|p|}_{\text{routing}} \;+\; \underbrace{O(1)}_{\text{iter-1 indicator}} \;=\; |p| + c_U,
\]
where $c_U := c_U' + O(1)$ collects all $p$-independent overhead.

\paragraph{Conclusion.} Choosing $p$ to be a shortest $U$-program for $s$, $|p| = K(s)$, gives $\nc(s)
  \le \norm{\theta_p}_0 \le K(s) + c_U$, proving (\ref{eq:upper}).

\section{Detailed Proof of the Lower Bound}\label{app:lower}

We give the self-delimiting encoding $\hat p : \Param_{\delta, M} \to \B^*$ of
Lemma~\ref{lem:lower}.

\paragraph{Step 1: enumerate non-zero parameters.} List the $W = \norm{\theta}_0$ non-zero parameters as $\Lambda(\theta) =
  \{(\ell_j, u_j, v_j, w_j)\}_{j=1}^{W}$, where $\ell_j \in \N$ is a layer index,
$u_j, v_j \in \N$ are source and target neuron indices, and $w_j \in
  \mathcal{V} := (\delta\Z \cap [-M, M]) \setminus \{0\}$ is the value. The value
alphabet $\mathcal{V}$ has size $2 \lfloor M/\delta \rfloor = O(1)$.

\paragraph{Step 2: prune and relabel.} Neurons not appearing in any $u_j, v_j$ contribute no computation and may be
removed; the network's output is unchanged. After pruning, the set of distinct
neuron indices has cardinality $n \le 2W$. Layers with no non-zero parameter
may be collapsed; the number of layers is at most $L' \le W$. Relabel neurons
and layers by order-preserving bijections to $\{1, \ldots, n\}$ and $\{1,
  \ldots, L'\}$.

\paragraph{Step 3: encode.} The description $\hat p(\theta)$ consists of:
\begin{itemize}
  \item A self-delimiting prefix encoding the integer $W$ (Elias-style): $\log_2 W + 2
          \log_2 \log_2 W + O(1) = O(\log W)$ bits.
  \item The state dimension $n$ and number of layers $L'$ in $O(\log W)$ bits each.
  \item Architecture metadata (head count, hidden width, activation choice) in $O(1)$
        bits.
  \item For each $j \in [W]$, the tuple $(\ell_j, u_j, v_j, w_j)$ in
        \[
          \lceil \log_2 L' \rceil + 2 \lceil \log_2 n \rceil + \lceil \log_2 |\mathcal{V}| \rceil \;\le\; 3 \log_2 W + O(1)
        \]
        bits, fixed-width binary; the widths are derivable from the prefix and so need
        not be transmitted.
\end{itemize}
The total length is $3 W \log_2 W + O(W)$ bits.

\paragraph{Step 4: simulator program.} Let $\Pi$ be a constant-size program for $U$ that:
\begin{enumerate}
  \item Parses the self-delimited prefix to recover $W, n, L'$, and the architecture
        metadata.
  \item Parses the $W$ tuples to reconstruct $\theta \in \Param_{\delta, M}$.
  \item Initialises $X^{(0)} = \mathbf{0}$ and iterates the network forward at fixed
        precision, applying each layer's linear update plus ReLU and the activation
        discretisation specified by the architecture.
  \item At each iteration $t \ge 1$, after each layer is applied: if $(X^{(t)})_{\halt}
          > 0$, halt; otherwise, if $(X^{(t)})_{\emit} > 0$, emit
        $\mathbf{1}[(X^{(t)})_{\bit} > 0]$.
\end{enumerate}
$\Pi$ is fixed, with constant size $|\Pi|$ independent of $\theta$.

\paragraph{Step 5: assembly.} The complete program is $\hat p(\theta) = \Pi \cdot \text{encoding}(\theta)$
(concatenation, parseable on $U$ since $\Pi$ knows where its own code ends).
Then $U(\hat p(\theta)) = s$ and
\[
  |\hat p(\theta)| \;\le\; |\Pi| + 3 W \log_2 W + O(W) \;\le\; c_d W \log_2 W + c_d
\]
for $c_d = \max(3 + \eps, |\Pi|)$ and all $W \ge 1$. Hence $K(s) \le c_d
  \norm{\theta}_0 \log \norm{\theta}_0 + c_d$; taking the infimum over $\theta$
outputting $s$ gives (\ref{eq:lower}).

\section{The Permutation Tightness Witness}\label{app:permutation}

We give the explicit construction underlying the permutation example in
\S\ref{sec:main:tightness}. The clean choice of output format is the one-hot
row-major serialisation of the permutation matrix; this avoids any
bit-extraction circuit and keeps the construction at $\Theta(N)$ non-zero
parameters.

\paragraph{Setup.} Fix $N \in \N$ and a permutation $\pi : [N] \to [N]$. Let $P_\pi \in \{0,1\}^{N
  \times N}$ denote its permutation matrix: $(P_\pi)_{i,\pi(i)} = 1$ and zeros
elsewhere. Define
\[
  s_\pi \;=\; (P_\pi)_{1,1}\,(P_\pi)_{1,2}\,\cdots\,(P_\pi)_{1,N}\,(P_\pi)_{2,1}\,\cdots\,(P_\pi)_{N,N} \;\in\; \B^{N^2},
\]
the row-major serialisation of $P_\pi$. The string $s_\pi$ has length $N^2$
with exactly $N$ ones and $N^2 - N$ zeros.

\paragraph{Kolmogorov complexity.} The map $\pi \mapsto s_\pi$ is a bijection, so describing $s_\pi$ is equivalent
to describing $\pi$. Counting gives $K(s_\pi) \ge \log_2 N! - O(1)$ for at
least half of all permutations, and the universal upper bound gives $K(s_\pi)
  \le \log_2 N! + O(\log N)$. Hence
\[
  K(s_\pi) \;=\; \log_2 N! + O(\log N) \;=\; \Theta(N \log N) \qquad \text{for typical (and worst-case) }\pi.
\]

\paragraph{Construction.} We exhibit a ternary looped network $\theta_\pi$ that outputs $s_\pi$ with
$\norm{\theta_\pi}_0 = \Theta(N)$. All $\pi$-dependence is concentrated in
exactly $N$ non-zero parameters; the remaining $\Theta(N)$ non-zero parameters
realise a fixed $\pi$-independent control circuit.

\emph{State coordinates.} The state has $n = \Theta(N)$ coordinates:
\begin{itemize}
  \item Coordinates $1, 2, 3$: halt, emit, bit channels.
  \item Coordinates $4, \ldots, 3 + N$: \emph{row indicator} $r \in \R^N$, a one-hot
        vector for the current row index $i \in [N]$.
  \item Coordinates $4 + N, \ldots, 3 + 2N$: \emph{column scanner} $c \in \R^N$, a
        one-hot vector for the current column index $j \in [N]$.
  \item One scratch coordinate $m$: the \emph{match indicator}, equal to
        $(P_\pi)_{i,j}$ at the current $(i, j)$.
\end{itemize}

\emph{Permutation parameters ($\pi$-dependent, exactly $N$ non-zeros).} For each $i \in [N]$, place a single $+1$ ternary parameter in the network's
``edge table'': specifically, a weight from row coordinate $r_i$ to a hidden
unit $h_i$, and the column scanner $c$ is wired into $h_i$ via a $\pi$-dependent
permutation map. The $N$ weights $W_\pi[k, \pi(k)] = +1$ encode the choice of
$\pi$; their positions carry $\log_2 N!$ bits of information.

\emph{Match circuit ($\pi$-independent, $\Theta(N)$ size).} A fixed FFN sub-layer computes the bilinear form $m = r^\top W_\pi c$. The
hidden layer has $N$ units, the $k$th computing $\text{ReLU}(r_k + (W_\pi c)_k -
  1)$ (an AND of two $\{0,1\}$ values). The output combines all hidden units with
$+1$ weights into the match channel. The fan-in from $r$ is the identity ($N$
non-zero diagonal weights); the fan-in from $c$ uses the $N$ permutation
parameters; biases ($-1$ per hidden unit) and the output projection ($+1$ per
hidden unit) contribute $\Theta(N)$ further non-zero parameters, all
$\pi$-independent.

\emph{Scanner advance ($\pi$-independent, $\Theta(N)$ size).} A fixed sub-circuit advances $c$ by one position per iteration via a cyclic-shift
matrix on $N$ coordinates ($N$ non-zero weights). When $c$ wraps from position
$N$ back to position $1$, the same circuit advances $r$ by one position
($\Theta(N)$ further non-zero weights). When $r$ wraps around, the halt channel
fires (constant-size logic).

\emph{Emission.} At each iteration, the emit channel is asserted (a constant
bias) and the bit channel is set to the match indicator $m$. After $N^2$
iterations the row-major serialisation of $P_\pi$ has been streamed to the
output and the halt channel fires.

\paragraph{Parameter count.}
\[
  \norm{\theta_\pi}_0 \;=\; \underbrace{N}_{\pi\text{-dependent}} \;+\; \underbrace{\Theta(N)}_{\pi\text{-independent control}} \;=\; \Theta(N).
\]
The $N$ $\pi$-dependent parameters encode the choice of permutation through
their positions, contributing $\log_2 N!$ bits of information. The
$\pi$-independent control circuit has fixed structure shared across all $\pi$
and all $N$, but instantiates with $\Theta(N)$ non-zero entries when
$N$-dimensional one-hot encodings are used.

\paragraph{Conclusion.} For typical (uniformly random) $\pi$, the construction shows
\[
  K(s_\pi) \;=\; \Theta(N \log N) \;=\; \Theta\!\bigl(\nc(s_\pi)\, \log \nc(s_\pi)\bigr),
\]
saturating (\ref{eq:lower}) up to constants. The logarithmic factor is
essential: the $N$ permutation parameters could not, by counting, identify
$\pi$ with fewer than $\log_2 N! = \Theta(N \log N)$ bits of positional
information, and that is exactly what the network exploits to emit the
$N^2$-bit string $s_\pi$. \qed

\paragraph{Remark on the binary-output variant.} If one prefers the more familiar output format $s'_\pi =
  \pi(1)\pi(2)\cdots\pi(N)$ written in binary (length $N \lceil \log_2 N \rceil$,
K-complexity still $\Theta(N \log N)$), the construction needs an additional
bit-extraction circuit that converts the one-hot column indicator into its
$\log_2 N$-bit binary representation. By the address-space argument of
\S\ref{sec:main:tightness}, any such fixed circuit on $N$-dimensional one-hot
vectors requires $\Theta(N \log N)$ non-zero parameters (equivalent to storing
a $\log_2 N \times N$ readout matrix), so the binary variant gives $\nc(s'_\pi)
  = \Theta(N \log N)$. The $L_p$ collapse and our main bound are unaffected, but
the binary variant does \emph{not} saturate the log factor. The one-hot variant
above is the cleaner witness.

\section{$L_p$ Sandwich Bounds}\label{app:lp}

\begin{proposition}[$L_p$ sandwich]\label{prop:lp}
  For every $p \in [1, \infty)$ and every computable $s$,
  \[
    \nc_p(s)^p \;\le\; M^p (K(s) + c_U) \quad \text{and} \quad K(s) \;\le\; \tfrac{c_d}{\delta^p} \nc_p(s)^p \log(\nc_p(s)^p / \delta^p) + c_d,
  \]
  where $\nc_p(s) = \inf_\theta \{ \norm{\theta}_p : \theta \in \Param_{\delta,
      M},\, \theta \text{ outputs } s \}$.
\end{proposition}

\begin{proof}
  By (\ref{eq:lp-collapse}), $\delta^p \norm{\theta}_0 \le \norm{\theta}_p^p \le M^p \norm{\theta}_0$, so $\nc_p(s)^p \le M^p \nc(s)$ and $\nc(s) \le \nc_p(s)^p / \delta^p$. Substitute into Theorem~\ref{thm:main}.
\end{proof}

Absorbing the precision constants:
\[
  L_1: \nc_1(s) = \tilde\Theta(K(s)), \qquad L_2: \nc_2(s)^2 = \tilde\Theta(K(s)), \qquad L_p: \nc_p(s)^p = \tilde\Theta(K(s)).
\]
Spectral and operator norms are not directly captured by parameter count:
spectral regularisation controls Lipschitz constants rather than information
content. The relationship between spectral norms and $K(s)$ is left open.

\section{Precision Regimes and the Limit of Vanishing Precision}\label{app:precision}

Theorem~\ref{thm:main} holds for any fixed precision; the constants depend on
the dynamic range $M/\delta$.

\paragraph{Ternary ($\delta = M = 1$).} Cleanest constants: per-entry description $3 \log_2 W + 1$ bits.

\paragraph{$b$-bit integer ($\delta = 1$, $M = 2^{b-1} - 1$).} Per-entry $3 \log_2 W + b$ bits. Linear in $b$.

\paragraph{$b$-bit dyadic.} Value alphabet of size $2^b$: per-entry $3 \log_2 W + b$ bits.

\paragraph{Quantitative dependence.} Both $c_U$ and $c_d$ scale as $O(\log_2(M/\delta))$. Even at very fine
precision, the sandwich bound holds with constants growing only as
$\log(M/\delta)$, i.e., as the bit-width.

\paragraph{The limit $\delta \to 0$.} The sandwich becomes vacuous: \citet{siegelmann1995computational} show some
real-weight networks compute non-recursive functions with finite norm; for such
networks $K(s) = \infty$ while $\norm{\theta}_2$ is finite. The fixed-precision
condition is what locks the result to the computable world. Since real hardware
always uses finite precision, the bound applies to every regime that runs on
actual deployments.

\section{Levin Complexity and Bounded Halting Time}\label{app:levin}

If we restrict halting time to $\le T$ iterations, the construction of
Lemma~\ref{lem:upper} translates a $U$-program $p$ with running time $T(p)$
into a network with halting time $T'(p) \le \rho \cdot T(p)$, where $\rho$ is
the per-step \emph{simulation overhead} of $T_U$: the number of looped-network
iterations per $U$-step. By the constructions of
\citet{li2025constant,giannou2023looped}, $\rho$ is a constant independent of
$p$, $|p|$, and the network state (it depends only on the architecture and on
the universal Turing machine $U$).

The encoding of Lemma~\ref{lem:lower} adds $\log T'$ bits to record the
iteration budget. Substituting $T' = \rho T$ and absorbing $\log \rho = O(1)$
into the constant gives
\[
  K(s) + \log T \;\le\; O\!\bigl(\nc_T(s) \log \nc_T(s)\bigr) + O(1),
\]
where $\nc_T(s)$ is the minimum non-zero parameter count of a network halting
within $T$ iterations and outputting $s$. The left-hand side is, up to
constants, Levin's $Kt$-complexity \citep{levin1973universal}, $Kt(s) = \min_p
  (|p| + \log T(p))$. Time-bounded neural complexity coincides with Levin
complexity up to log factors, with the multiplicative simulation overhead
$\rho$ of $T_U$ contributing only an additive $\log \rho = O(1)$ to the bound.
This form is most directly relevant to learning, since training procedures cap
the running time of the learned model.

\section{Proof of the Solomonoff Connection}\label{app:prior}

\paragraph{The two priors and their equivalence.} The Gaussian prior over fixed-precision weights, $\pi_{L_2}(\theta) \propto
  2^{-\norm{\theta}_2^2}$, and the sparsity prior, $\pi_{L_0}(\theta) \propto
  2^{-\norm{\theta}_0}$, agree up to constants in the exponent: by
(\ref{eq:lp-collapse}),
\[
  2^{-M^2 \norm{\theta}_0} \;\le\; 2^{-\norm{\theta}_2^2} \;\le\; 2^{-\delta^2 \norm{\theta}_0}
  \qquad \text{for } \theta \in \Param_{\delta, M},
\]
so $\pi_{L_2}(\theta) = \pi_{L_0}(\theta)^{\beta(\theta)}$ for some
$\beta(\theta) \in [\delta^2, M^2]$. The same collapse gives the analogous
relation between $\pi_{L_2}$ and $\pi_{L_1}(\theta) \propto
  2^{-\norm{\theta}_1}$. The calculation below uses $\pi_{L_0}$ for clarity; the
result transfers to any $L_p$ weight decay by absorbing the precision constants
into $\alpha, \beta$. Let $Q(s) = \sum_{\theta \,:\, \theta \text{ outputs } s}
  \pi_{L_0}(\theta)$.

\paragraph{Uniformity of constants in $s$.} The constants $\alpha, \beta$ in Corollary~\ref{cor:prior} depend on $(\delta,
  M)$ and the architecture but not on $s$. This is automatic from the proof: the
upper bound uses $\norm{\theta_*}_0 \le K(s) + c_U$ with $c_U$ uniform; the
lower bound uses $\norm{\theta}_0 \ge K(s)/(c_d \log K(s)) - O(1)$ with $c_d$
uniform, plus the geometric-series argument whose multiplicative slack depends
only on the counting bound $W^{O(W)}$ and not on $s$.

\paragraph{Lower bound on $Q(s)$.} By Lemma~\ref{lem:upper}, some $\theta_*$ outputs $s$ with $\norm{\theta_*}_0
  \le K(s) + c_U$. So
\[
  Q(s) \;\ge\; \pi(\theta_*) \;=\; 2^{-\norm{\theta_*}_0} \;\ge\; 2^{-K(s) - c_U}.
\]

\paragraph{Upper bound on $Q(s)$.} By Lemma~\ref{lem:lower}, every $\theta$ outputting $s$ has $\norm{\theta}_0
  \ge K(s)/(c_d \log K(s)) - O(1)$. The number of distinct $\theta \in
  \Param_{\delta, M}$ with $\norm{\theta}_0 = W$ is at most $W^{O(W)}$ (counting
(location, value) tuples). Summing $W^{O(W)} \cdot 2^{-W}$ over $W \ge
  K(s)/(c_d \log K(s))$ gives a geometric series dominated by its first term:
\[
  Q(s) \;\le\; 2^{-K(s)/(c_d \log K(s)) + O(\log K(s))} \;=\; 2^{-K(s)/(c_d' \log K(s))}.
\]

\paragraph{Comparison with Solomonoff.} The universal prior $M$ satisfies $-\log M(s) = K(s) + O(1)$
\citep{li2008introduction}. The L2-weight-decay prior $Q$ matches in the
exponent up to the logarithmic factor of our bound: $-\log Q(s) \in [K(s)/(c_d'
    \log K(s)),\, K(s) + c_U]$.

\section{Computability of $\nc$}\label{app:computability}

\begin{proposition}
  $\nc$ is upper-semicomputable but not computable.
\end{proposition}

\begin{proof}
  \emph{Upper-semicomputability.} Enumerate fixed-precision looped networks in increasing $\norm{\cdot}_0$ order, simulate each on $x_0 = 0$ with progressively larger time bounds (dovetailing), and emit a new upper bound on $\nc(s)$ each time a network halts and outputs $s$.

  \emph{Non-computability.} By Theorem~\ref{thm:main}, $\nc(s) = \tilde\Theta(K(s))$. If $\nc$ were computable, $K$ would be approximable to within a logarithmic factor by a computable function, contradicting the standard non-computability of $K$ \citep{li2008introduction}.
\end{proof}

%%%%%%%%%%%%%%%%%%%%%%%%%%%%%%%%%%%%%%%%%%%%%%%%%%%%%%%%%%%%

% \newpage
% \input{checklist.tex}

\end{document}
